\section{Related Work}
\label{sec:resources}

Large language models (LLMs) have been investigated as translation systems through studies of
translation quality, prompting, and adaptation. Hendy et al.\ compare GPT models with research
and commercial systems, reporting competitive high-resource performance but limitations for
low-resource languages in their evaluation \cite{hendy2023}. Vilar et al.\ demonstrate the
importance of demonstration quality for prompted translation \cite{vilar2023}, while Moslem
et al.\ investigate adaptation using in-domain examples and terminology \cite{moslem2023}.
These studies motivate assessing translation capability under explicitly defined conditioning
rather than assuming uniform performance across languages and domains.

Translation benchmarks operationalize these comparisons. FLORES-101 provides professionally
translated, fully aligned multilingual references for many-to-many evaluation \cite{goyal2022}.
WMT25 compares submitted systems, LLMs, and online providers using challenging test material,
multiple domains, and document context \cite{kocmi2025}. Its professional evaluation covers
15 of the 30 offered language pairs using Error Span Annotation or Multidimensional Quality
Metrics (MQM). This work underscores the importance of test difficulty, language coverage,
and evaluation context; our segment-level benchmark does not replicate its document-level protocol.

Model comparison requires complementary quality measures. BLEU measures reference $n$-gram
overlap \cite{papineni2002}, SacreBLEU standardizes its reporting \cite{post2018}, and chrF++
combines character and word overlap \cite{popovic2017}. COMET learns quality prediction from
human judgments \cite{rei2020}, whereas BERTScore compares contextual token representations
\cite{zhang2020}. Expert MQM evaluation demonstrates the importance of annotator expertise
and document context \cite{freitag2021}. Terminology accuracy supplies an additional dimension,
as reflected in dedicated WMT evaluation \cite{alam2021}; variation-aware analysis further
distinguishes legitimate alternatives from undesirable inconsistency \cite{dahan2026}.
Reference-term retention, however, is narrower than bilingual terminology accuracy: span presence
does not establish contextual correctness, and reference mismatch need not indicate mistranslation.

Terminology extraction supports the annotations needed for these diagnostics. Chun et al.\ use
syntactically retrieved demonstrations to improve LLM term-boundary prediction on specialized
benchmarks \cite{chun2025}. Such work motivates evaluating extraction quality independently:
verifying occurrence in a reference does not establish annotation completeness or expert validity.

In chemistry, related tasks must be distinguished from prose translation. STOUT studies molecular
representation-to-name conversion with structure-oriented checks \cite{rajan2021}, while CHEMDNER
evaluates chemical mention recognition against expert annotations \cite{krallinger2015}; neither
benchmarks multilingual chemical prose translation. For legal comparison, JRC-Acquis
\cite{steinberger2006}, distributed through OPUS \cite{tiedemann2012}, provides parallel text.
LegISTyr distinguishes terminology insertion from contextual choice and fluency \cite{dinatale2025},
and an EEA terminology suite reports performance that need not follow general MT rankings
\cite{hauksdottir2025}. Our study centers on multilingual chemical-patent translation capability,
using legal text as a comparison domain and terminology metrics as complementary diagnostics.
Reference-derived terminology hints constitute a separate oracle-assisted condition, not evidence
of source-only capability.